%=======================================================================
%  APPENDICES
%=======================================================================
\section*{Appendix A -- AI Use Statement}
\addcontentsline{toc}{section}{Appendix A -- AI Use Statement}
\setcounter{table}{0}
\renewcommand{\thetable}{A.\arabic{table}}

Relevant AI-supported activities are documented below. Not every prompt is
listed; the focus is on how AI contributed to the work and how important
outputs were verified. AI output is not used as evidence for factual claims.

\begin{longtable}{L{2.4cm} L{3.0cm} L{4.6cm} L{4.2cm}}
\caption{AI use statement.}\label{tab:aiuse}\\
\toprule
\textbf{AI tool} & \textbf{Used for} & \textbf{How the output influenced the project} & \textbf{How important output was verified} \\
\midrule
\endfirsthead
\toprule
\textbf{AI tool} & \textbf{Used for} & \textbf{How the output influenced the project} & \textbf{How important output was verified} \\
\midrule
\endhead
\bottomrule
\endfoot

AskNature (search and strategy pages) & Discovering biological strategies for functions F1--F12, navigated by taxonomy branch rather than by organism &
Produced the candidate list in Table~\ref{tab:models}, including four models that were subsequently de-selected &
Every mechanism claim carried into the report is referenced. The pages could not be opened programmatically, so titles and URLs were taken from search results; the open item in Section~\ref{sec:discovering} records that each must be confirmed and, where a primary paper exists, replaced by it \\

Large language model assistant & Structuring the scoping, distilling the course material into working notes, drafting report text, building the transient wall model of Section~\ref{sec:creating}, and generating prompts for figures &
Draft text and table structures throughout Sections~1--4; the search-path table; the transient model and its figures &
No factual claim entered the report without a named source. Benchmark values were checked against the Fachplanung. Two errors in the AI-built wall model were found and corrected: an air-cavity node made the explicit scheme unstable, and the synthetic weather drive had an inverted diurnal phase, which showed up as a peak surface temperature exactly equal to the air temperature \\

AI as critical reviewer & \todo{Section~3.3 asks for this. Give the reviewer the concept sketch, the Engineering Canvas, the functions and KPIs, and the assumptions} &
\todo{Record each comment and the team's decision: accept, reject or investigate} &
\todo{A documented \emph{reject} with a reason is worth more here than a list of accepted comments, because it shows judgement} \\

\todo{Further tools used by the team} & \todo{} & \todo{} & \todo{} \\
\end{longtable}

\paragraph{On the limits of this statement.} Nothing in this report is cited to
an AI tool. Where an AI-assisted search surfaced a candidate, the candidate was
carried forward only with a source of its own, and the two remaining places where
that chain is incomplete are marked as open items rather than left implicit: the
AskNature pages listed above, and the \qty{23}{\percent} figure for elephant
hair, which traces to a university blog rather than to a study
\parencite{notredame_elephant}.

\section*{Appendix B -- Decision Log}
\addcontentsline{toc}{section}{Appendix B -- Decision Log}

Every decision below is argued at the place named in the last column. This table
exists so that the reasoning can be checked without re-reading the report, and
so that what was \emph{not} chosen stays visible: a design process is only
legible if the discarded alternatives are on the record.

\renewcommand{\thetable}{B.\arabic{table}}
\setcounter{table}{0}

\begin{longtable}{L{2.6cm} L{3.0cm} L{3.0cm} L{5.7cm}}
\caption{Decision log. What was chosen, what it displaced, and why.}
\label{tab:decisionlog}\\
\toprule
\textbf{Decision} & \textbf{Chosen} & \textbf{Rejected alternative} & \textbf{Reason, and where it is argued} \\
\midrule
\endfirsthead
\toprule
\textbf{Decision} & \textbf{Chosen} & \textbf{Rejected alternative} & \textbf{Reason, and where it is argued} \\
\midrule
\endhead
\bottomrule
\endfoot

Design question & Let the absorbed heat drive its own removal, self-regulating & Four other candidates, from \enquote{prevent heat build-up in cities} to \enquote{turn heat into a resource} & Too broad, unanchored, factually wrong, or not verifiable in the time available. The chosen question adds the self-regulating clause, which is the actual content of the project. Section~\ref{sec:designquestion} \\

Reference city & Zurich & A generic European city & An adopted municipal heat strategy, published target values, open geodata and long-term stations make the targets checkable instead of invented. Section~\ref{sec:context} \\

Reference site & Modellierungsgebiet 03, closed perimeter block & A freely chosen building & The city publishes modelled baseline and climate-optimised results for exactly this area, so the concept is compared against a documented reference. Section~\ref{sec:context} \\

Design case & Five consecutive hot days & One hot afternoon & Only a multi-day case stresses the water reserve and the reset behaviour of the store; a concept that works on day one and saturates by day three has not solved the problem. Section~\ref{sec:context} \\

Evening window & 18:00--24:00 & 20:00--24:00 & A window must contain the instant at which its own KPI is evaluated, and the two hours before 20:00 are among the hottest of the day at the surface. Section~\ref{sec:context} \\

System of interest & The façade, with the roof as the release interface & The envelope as a whole, roof included & The roof receives more than twice the noon irradiance, but releases its heat upward to a sky it can see; the façade releases sideways into the canyon at head height, and Section~1 fixes that as where the effect must appear. Section~\ref{sec:context} \\

Number of systems & One & Two --- \enquote{reduce heat} and \enquote{prevent heat} & The systems view showed these are two functions of one system, coupled by the design question. Section~\ref{sec:systemview} \\

Scope of \enquote{use the heat} & Let it drive its own removal & Seasonal storage and electricity generation & Conversion competes with photovoltaics for the same surface and any seasonal store lies outside the envelope. Section~\ref{sec:systemview} \\

Selection KPI for the urban effect & PET at \qty{2}{\metre} above ground & Urban heat island intensity & A district-scale quantity that a single retrofit cannot move measurably, and it invites the surface-versus-air confusion. Section~\ref{sec:kpi} \\

Cost and embodied emissions & Constraints & An eighth KPI & A KPI must follow from a function; these are conditions to satisfy, not measures of performance. The move also brings the list inside the required four to seven. Section~\ref{sec:kpi} \\

Albedo target & $\geq0.80$, the city's threshold & The technical maximum of $\approx0.98$ & Life's Principles ask for optimised rather than maximised; further reflectance costs glare, winter gains and photovoltaic yield. Section~\ref{sec:lifesprinciples} \\

Life's Principles claimed & Four claimed with the sub-principle named & \enquote{Creates opportunities for life} and \enquote{Evolve to survive} for the artefact & A damp textured surface invites the algal growth KPI~7 penalises, and a façade layer does not evolve across generations. Naming the conflict is more honest than claiming the principle. Section~\ref{sec:lifesprinciples} \\

Search strategy & By function, through the taxonomy & By organism & Searching for organisms would not have produced F9, which is the pairing of two taxonomy branches rather than a single one. Section~\ref{sec:researchpath} \\

Termite mound & De-selected & The standard biomimicry reference for passive cooling & The chimney account is contested in the literature, and the building usually cited for it uses fans. Section~\ref{sec:discovering} \\

Fog harvesting & De-selected & The Namib beetle, the most-cited water model & Fog in Zurich is a winter phenomenon; a heat wave brings neither rain nor fog. Section~\ref{sec:discovering} \\

Green façades & Out of scope & An obvious and proven measure & Already the client's HA~09 and HA~10, and in direct competition with our own water reserve, which is scarcest exactly when planting needs most. Section~\ref{sec:discovering} \\

Evidence standard & Primary papers where they exist, the publishing organisation otherwise & AskNature entries treated as evidence in their own right & AskNature and AI-assisted search find candidates. Twelve of twenty model entries still rest on an AskNature page alone; this is recorded as a shortfall rather than presented as met. Section~\ref{sec:researchpath} \\

Concept pair & \todo{decide} & \todo{the two directions not developed} & \todo{state the physical distinction that makes the pair meaningful} \\

Final concept & \todo{A, B or a justified hybrid} & \todo{} & \todo{name the weakness in V1 that V2 answers} \\

\end{longtable}

\clearpage
\section*{Appendix C -- Scoping Detail}
\addcontentsline{toc}{section}{Appendix C -- Scoping Detail}
\setcounter{table}{0}
\renewcommand{\thetable}{C.\arabic{table}}

\noindent Three tables that document the scoping decisions of
Section~\ref{sec:designquestion} to Section~\ref{sec:systemview}. They are
placed here rather than in the body so that Section~2 stays within the page
budget while the reasoning behind it remains on record.

\begin{table}[htbp]
\centering
\caption{Design questions considered, and why four were rejected.
Referenced from Section~\ref{sec:designquestion}.}
\label{tab:questions}
\renewcommand{\arraystretch}{1.2}
\small
\begin{tabularx}{\textwidth}{L{0.5cm} X L{4.6cm}}
\toprule
\textbf{\#} & \textbf{Candidate} & \textbf{Assessment} \\
\midrule
1 & How do we prevent heat build-up in urban areas?
  & \emph{Too broad.} No place, no time, no intervention point --- the
    equivalent of \enquote{How might we end hunger?} \\
2 & How could we use the energy or heat build-up positively while lowering
    temperature?
  & Right idea, but unanchored: no location, and \enquote{positively} does not
    say for what. Refined into candidate~5. \\
3 & How might we get the heat that Zurich's existing buildings store during the
    day back out of the city before the next morning, without using energy?
  & Well formed and highly testable. Rejected only because it discards the
    team's central idea of \emph{using} the heat. Its night-time focus is
    retained in the selected question. \\
4 & How can we prevent heat waves causing heat islands?
  & \emph{Factually wrong.} Heat waves do not cause heat islands. The urban
    heat island is a permanent consequence of urban surface properties and
    geometry and exists on ordinary summer days; a heat wave makes it
    dangerous, it does not create it \parencite{oke1982}. \\
5 & How might we turn the heat that accumulates in a city into a useful
    resource, instead of reflecting or discarding it?
  & Good direction, but the reading \enquote{store summer heat for winter} is
    not verifiable within the project, and it has no biological support:
    organisms store chemical energy across seasons, not sensible heat.
    Narrowed into the selected question. \\
\bottomrule
\end{tabularx}
\end{table}

\begin{table}[htbp]
\centering
\caption{Context conditions and constraints in full. Hard constraints
disqualify a concept when violated; soft constraints weight the evaluation in
Section~\ref{sec:evaluation}. Summarised in Section~\ref{sec:context}.}
\label{tab:constraints}
\renewcommand{\arraystretch}{1.2}
\small
\begin{tabularx}{\textwidth}{L{2.3cm} L{1.4cm} X}
\toprule
\textbf{Type} & \textbf{Class} & \textbf{Constraint} \\
\midrule
Energy & hard & No auxiliary electrical energy for the cooling function itself. This follows directly from the design question: the heat load must drive the response. Sensors used for validation are excluded. \\
Water & hard & No potable water. Only rainwater harvested on site or atmospheric moisture. \\
Existing building & hard & Retrofit onto an existing envelope. Load-bearing structure, floor plan and window openings are not redesigned. \\
Townscape & hard & The result must be acceptable in the Zurich townscape. Buildings may be listed; heritage protection, building law and neighbour rights can restrict visible changes \parencite[128]{stadtzuerich2020}. \\
Glare & hard & No specular glare towards streets, neighbouring flats or traffic. The city names glare as the main objection to bright envelopes \parencite[128]{stadtzuerich2020}, which favours diffuse over mirror-like surfaces. \\
Climate & hard & Must survive the Swiss annual cycle: frost to about \qty{-10}{\celsius}, freeze--thaw cycles, driving rain, UV, and periods of several dry weeks in summer. \\
Winter & hard & Must not appreciably increase heating demand. A high solar reflectance also reduces useful winter solar gains; this trade-off must be quantified, not ignored. \\
Installation & soft & Must be installable by ordinary trades on an occupied building, ideally without full scaffolding. \todo{state an installation-time target per \unit{\square\metre}} \\
Maintenance & soft & Must remain functional with little maintenance. Soiling and the resulting cleaning cost are named by the city as a specific weakness of high-albedo envelopes \parencite[128]{stadtzuerich2020}. Service life target $\geq$ \qty{20}{years}. \\
Solar energy & soft & Should not unnecessarily compromise roof- or façade-integrated photovoltaics; the city notes that bright envelopes reduce PV yield \parencite[128]{stadtzuerich2020}. State which envelope areas the concept claims. \\
Structural & soft & Limited additional dead load and build-up thickness on an existing envelope. \\
Embodied emissions & soft & The greenhouse gas embodied in the added build-up must be repaid by the avoided cooling demand within the service life. Measured in \unit{\kilo\gram}\,CO$_2$-eq\,\unit{\per\square\metre}. \todo{source a benchmark} \\
Economic & soft & The city expects higher cost for high-albedo envelopes \parencite[128]{stadtzuerich2020}. Target: cost per square metre in the range of a normal refurbishment. \todo{find a benchmark figure and cite it} \\
\bottomrule
\end{tabularx}
\end{table}

\begin{table}[htbp]
\centering
\caption{What the comparison with each parallel system contributed to the
scoping. Referenced from Figure~\ref{fig:systemview} and
Section~\ref{sec:systemview}.}
\label{tab:parallel}
\renewcommand{\arraystretch}{1.2}
\small
\begin{tabularx}{\textwidth}{L{3.6cm} X}
\toprule
\textbf{Parallel system} & \textbf{What the comparison adds} \\
\midrule
Stormwater and sewer system
& Rainwater is currently drained away as fast as possible. Seen from our system
  it is a cooling reserve that is being discarded --- and during a heat wave it
  is the only water available under the \enquote{no potable water} constraint.
  This adds the functions \emph{capture}, \emph{retain} and \emph{meter} water
  (F6--F8) and makes storage capacity a design variable rather than an
  afterthought. \\
Street and pavement network
& The street is treated by the same client with the same logic (HA~06,
  Table~\ref{tab:incumbent}), and it achieves a slightly larger daytime effect
  than the envelope. The comparison sets a benchmark we must beat, and it shows
  that the envelope's advantage has to come from the night window, where the
  street performs no better. \\
Power grid and building energy system
& The obvious way to \enquote{use} heat is to convert it. The comparison shows
  why that is a different project: conversion competes with photovoltaics for
  the same surface, and any seasonal store lies outside the envelope. It is
  therefore placed out of the boundary, while the milder reading of
  \enquote{use} --- letting the heat drive its own removal --- stays in. \\
Residents, transport and public transit
& Defines where the effect has to be delivered: at pedestrian level in the
  street canyon and in the flats on the upper floors, not on the roof membrane.
  This is why PET at \qty{2}{\metre} above ground becomes a KPI, and it fixes
  the evening window 18:00--24:00, when people are outdoors and then sleeping. \\
\bottomrule
\end{tabularx}
\end{table}
